\section{Introduction}
\label{sec:introduction}

A convergence theorem describes the limiting behavior of an infinite
sequence, whereas a numerical method must stop at a finite iterate and decide
what that iterate certifies.  The choice of stopping quantity is therefore
part of the mathematical meaning of an algorithm.  The objective error
\(f(x)-f^\star\), where \(f^\star\) is the optimal value, answers the desired
question directly but is not generally computable.  A
Karush--Kuhn--Tucker (KKT) residual is usually
computable, but its interpretation changes with the geometry of the problem:
in convex optimization it is an optimality residual, while in nonconvex
optimization it is only a stationarity residual.

For a convex--concave saddle problem, the primal--dual gap avoids this
ambiguity.  It is nonnegative and vanishes exactly at a saddle point, thereby
combining primal optimality, dual optimality, and feasibility in one scalar.
The classical gap is nevertheless poorly suited to affine equality
constraints on unbounded spaces.  The supremum over the multiplier is
infinite whenever the current point is infeasible, and the primal comparison
may also be infinite.  Penalizing both comparison variables quadratically
leads to the \emph{smoothed duality gap} (SDG).  Its regularized comparison
problems are proximal problems, so the gap can remain finite away from
feasibility and can be evaluated without knowing an optimal solution.
Smooth primal--dual gap constructions were developed for convex composite
minimization in \cite{TranDinhFercoqCevher2018}; the self-centered SDG was
introduced in \cite{Fercoq2023} and then studied systematically as a stopping
criterion in \cite{WalwilFercoq2025}.

The convex interpretation cannot be transferred merely by removing the word
``convex''.  If \(f\) is proper, lower semicontinuous, and convex, every KKT
pair is a saddle point of the ordinary Lagrangian and yields attained primal
and dual optima with zero duality gap.  Conversely, when a primal solution
exists and the affine feasible set intersects
\(\operatorname{ri}(\operatorname{dom}f)\), strong duality and dual attainment
hold and the solution admits a multiplier.
Under these explicit hypotheses, the convex theory identifies
\[
  \text{saddle point}
  \quad\Longleftrightarrow\quad
  \text{KKT pair}
  \quad\Longleftrightarrow\quad
  \text{zero-gap primal--dual optimum}
\]
see \citet[Sections~5.4.2 and~5.5.3]{BoydVandenberghe2004} for the smooth
formulation and \citet[Theorem~27.4]{Rockafellar1970} for the convex
subdifferential condition under relative-interior qualification.  These
equivalences break down in a nonconvex problem.  A KKT pair
may represent a local minimum, a saddle, or a local
maximum.  Consequently, before using the SDG beyond convexity, one must
identify its exact zero set and distinguish the statement ``the gap is zero''
from the much stronger statement ``the point is globally optimal''.  A
zero-based stopping rule with the wrong interpretation would accept every KKT
branch, including saddles and local maxima.

\subsection{Problem setting}
\label{subsec:problem-setting}

Let \(A\in\R^{m\times n}\) and \(b\in\R^m\), and define the affine feasible
set
\[
  \Cset:=\{x\in\R^n:Ax=b\}.
\]
No full-row-rank assumption is imposed on \(A\), and we assume that
\(\Cset\cap\operatorname{dom}f\ne\varnothing\).  We study
\begin{equation}
  \min_{x\in\R^n} f(x)
  \quad\text{subject to}\quad
  Ax=b,
  \label{eq:main-problem}
\end{equation}
where \(f\) is proper, lower semicontinuous, and \(\rho\)-weakly convex for a
valid global modulus \(\rho\ge0\).  The chosen \(\rho\) need not be the
smallest possible modulus.  By definition,
\(f+\frac{\rho}{2}\norm{\cdot}^{2}\) is convex: the negative curvature of
\(f\) is globally bounded by a quadratic model.  This class includes convex
functions, smooth objectives with a uniform lower Hessian bound, and many
nonsmooth nonconvex models used in learning and signal recovery
\cite{DavisDrusvyatskiy2019}.  Controlled-concavity penalties such as SCAD
\cite{FanLi2001} and MCP \cite{Zhang2010} are standard examples.  Weak
convexity is therefore a useful intermediate regime: it is broad enough to
contain genuine nonconvexity, yet a proximal term with curvature larger than
\(\rho\) makes the regularized subproblem strongly convex.

Affine constraints retain the essential primal--dual difficulty while
keeping the geometry transparent.  They introduce feasibility, multipliers,
and tangent and normal directions, but they also permit an exact nullspace
description of feasible perturbations and a dual-root representation of the
constrained proximal map.  Problem~\eqref{eq:main-problem} is consequently a
controlled setting in which both the meaning of the SDG and the cost of
evaluating the proximal objects behind it can be studied.  We write its
ordinary Lagrangian as
\(\Lagr(x,y):=f(x)+\ip{y}{Ax-b}\), with multiplier \(y\in\R^m\).

The proximal curvature creates two distinct regimes.  If the primal
smoothing curvature \(\beta_x\) is larger than \(\rho\), the inner SDG
problem is strongly convex and has a unique global minimizer; we call this
the uniform safe regime.  Below this threshold, the inner problem may be
nonconvex and a local routine is not an oracle for the value-defined gap.
At the same time, the global modulus
\(\rho\) can be conservative at a fixed KKT pair.  The least curvature that
supports the Lagrangian globally at that pair is a pointwise quantity, and
different KKT pairs of the same problem may have different thresholds.  This
distinction between a uniform safe regime and pointwise global support is
closely related to the broader algorithmic difficulty of choosing a proximal
parameter when the weak-convexity modulus is unknown
\cite{Xie2026}.

The available certificates differ both in conclusion and in evaluation cost.
A KKT residual is local and computable.  In the safe regime, the SDG uses the
Lagrangian proximal point
\[
  p_{\beta_x}(x,y)
  :=\argmin_u\left\{\Lagr(u,y)
       +\frac{\beta_x}{2}\norm{u-x}^2\right\}
\]
together with an explicit feasibility term, and its zero set is the KKT set.
The rejection stage uses a different map, the hard-constrained safe prox
with \(\lambda>0\) and \(\lambda\rho<1\),
\[
  P_\lambda^{\Cset}(z)
  :=\argmin_{u\in\Cset}\left\{f(u)
       +\frac{1}{2\lambda}\norm{u-z}^2\right\}.
\]
A zero global augmented-Lagrangian gap instead identifies a primal global
minimizer and an augmented-dual maximizer
\cite{Rockafellar1974,YalcinKasimbeyli2020}, but requires a global infimum of
a generally nonconvex augmented Lagrangian and an attaining multiplier.  This
leads to the intermediate question whether the hard-constrained safe prox can
provide one-sided information beyond KKT without solving that global problem.

The analysis is organized around three precise questions.  First, which
variational condition is equivalent to zero SDG at arbitrary positive
smoothing, and how does it simplify in the uniform regime
\(\beta_x>\rho\)?  This question also asks for the least curvature required
at one fixed KKT pair.  Second, once a stationary candidate has been found,
can the hard-constrained safe proximal map exhibit a feasible point with
strictly lower objective value, including when a reduced-Hessian test is
inconclusive, and at what perturbation scale?  Third, on problems with
complete analytic stationary catalogs, how do the SDG, Lagrangian proximal
fixed-point residual, and KKT residual recover KKT branches, and how does
proximal escape classify the recovered candidates?  These questions
separate three notions that coincide in convex theory:
\[
  \text{first-order stationarity},\qquad
  \text{local minimality},\qquad
  \text{global optimality}.
\]

\subsection{Main results}
\label{subsec:main-results}

The first main result gives an exact description of what the weakly convex
SDG certifies.  In the uniform safe regime, its zero set coincides with the KKT
set, and the gap controls feasibility and a Lagrangian proximal fixed-point
residual separately.  For an arbitrary positive curvature, zero gap is characterized
through the level-proximal subdifferential.  This yields the pointwise
threshold
\[
  \beta_{\crit}(x,y)
  =
  \sup_{\substack{u\in\dom f\\u\ne x}}
  \frac{2[\Lagr(x,y)-\Lagr(u,y)]_+}{\norm{u-x}^{2}},
\]
which is bounded above by \(\rho\) at a KKT pair but may be strictly smaller.
Exact quartic examples show both pair-dependent thresholds and the false
zeros produced by a local inner solve.

The second main result is a post-KKT rejection framework based on the
proximal-escape radius, defined as the infimum of center-perturbation radii
admitting a lower proximal image.  The ideal radius separates global minima,
points that are not local minimizers, and nonglobal local minima.  A
constrained Moreau-envelope surrogate gives a smooth outer search, while
tangent-space reduction and an affine dual root implement the inner prox
without replacing feasibility by a finite penalty.  Because the argument uses
realized objective decrease rather than a quadratic model, it can also reject
degenerate stationary points at which the reduced Hessian is inconclusive.

The numerical study uses complete analytic stationary sets to compare branch
discovery with post-KKT rejection, including a controlled shrinking-radius
test at Hessian-flat stationary points.

The work reported here comprises
the analysis in \cref{sec:safe-sdg,sec:thresholds}, the escape classification
and affine constructions in \cref{sec:escape}, and the analytic and numerical
studies in \cref{sec:experiments}; the variational tools and connections to
existing approaches are cited in \cref{sec:preliminaries,sec:related-work}.

\subsection{Related work}
\label{sec:related-work}

The SDG analysis builds on the convex theory of smoothed gaps.
Tran-Dinh, Fercoq, and Cevher introduced a smooth primal--dual framework for
composite convex minimization \cite{TranDinhFercoqCevher2018}.  Fercoq
introduced the self-centered smoothed gap \cite{Fercoq2023}; Walwil and Fercoq
then developed it as a stopping criterion and compared it with the KKT error,
projected duality gap, and objective-gap/feasibility measure
\cite{WalwilFercoq2025}.  Their zero-gap interpretation uses the
convex equivalence between KKT conditions, saddle points, and primal--dual
optimality.  Here the same ordinary Lagrangian and self-centered smoothing are
examined after this equivalence fails.

Weakly convex Moreau-envelope theory provides regularity for both proximal
constructions above, although their objectives and feasible domains differ.
In the safe regime the envelope is continuously differentiable and its
gradient is a standard stationarity measure \cite{DavisDrusvyatskiy2019}.
Fixed points of the proximal map characterize stationarity rather than
global optimality.  Xie's
adaptive proximal framework addresses the
complementary problem of taking accepted steps when the global
weak-convexity parameter is unknown \cite{Xie2026}; the pointwise threshold
here instead asks whether one fixed center globally minimizes its own
proximal model.

The variational object connecting the SDG zero set with global proximal
support is the level-proximal subdifferential.  It represents proximal maps of
possibly nonconvex functions as resolvents without restricting the function
to the hypoconvex case \cite{WangWang2024}.  Unlike a local proximal
subdifferential, it imposes a quadratic support inequality against every
comparison point, precisely matching the global inner infimum in the SDG.
This supplies the pointwise language for the zero-set and threshold results.

Negative reduced-Hessian curvature is the classical local signature of a
strict saddle, but it is silent at degenerate stationary points whose first
and second derivatives vanish.  Higher-order escape methods use Taylor
information beyond the Hessian \cite{AnandkumarGe2016}, while linear
regularization can restore strict-saddle geometry in selected models
\cite{UblHaleYazdani2023}.  For weakly convex objectives,
\citet{DavisDrusvyatskiy2022} establish avoidance of active strict saddles
by proximal methods under suitable structural and initialization assumptions.
The present escape test asks a different question: whether proximal
polishing at perturbed centers supplies a lower feasible value for a fixed
candidate, including when negative Hessian curvature is absent.

For a stronger global comparison, Rockafellar relates augmented Lagrangians
to perturbation and stability
\cite{Rockafellar1974}, while Din\c{c} Yal\c{c}\i n and Kasimbeyli develop zero-gap
conditions through weak conjugacy and sharp augmentation
\cite{YalcinKasimbeyli2020}.  Their global conclusions require the relevant
stability, duality, and attainment hypotheses, and their dual functions retain
a global inner optimization problem.  Weak conjugacy is an
abstract-convexity construction, distinct from the \(\rho\)-weak convexity
used here to regularize proximal subproblems.

The algorithmic literature instead supplies appropriate comparisons for
finding constrained stationary points.  MEAL/LiMEAL
\cite{ZengYinZhou2022}, proximal augmented-Lagrangian methods
\cite{ZhangLuo2020}, feasible descent and SQP schemes
\cite{Rosen1960,LaleeNocedalPlantenga1998}, and distributed
Douglas--Rachford splitting for weakly convex multi-block models
\cite{HuXieCaiHan2025} are designed to approach stationary points or local
solutions.  Efficient computation of affine-constrained proximal operators is
itself a separate numerical problem; dual reduction and semismooth techniques
for structured nonconvex prox subproblems are developed in
\cite{HouTangTohWang2026}.  These works provide algorithmic context for
stationary-point computation, while the constrained proximal constructions
below support post-KKT rejection.

The remainder is organized as follows.
\Cref{sec:preliminaries} introduces the KKT system, weak convexity and
level-proximal support, followed by proximal maps, Moreau envelopes, and the
SDG.
\Cref{sec:safe-sdg} establishes the safe zero-set theorem.
\Cref{sec:thresholds} gives the pointwise support characterization and the
counterexamples beyond the uniform threshold.  \Cref{sec:escape} develops
the proximal-escape radius and its affine-constrained computation.
\Cref{sec:experiments} reports the numerical experiments, and
\Cref{sec:discussion} discusses their scope and concludes the report.
